\documentclass[12pt]{article}

\usepackage{amsmath,amssymb,amsthm,mathtools}
\usepackage[margin=1in]{geometry}
\usepackage{enumitem}
\usepackage{xurl}
\usepackage[hidelinks]{hyperref}

\theoremstyle{plain}
\newtheorem{theorem}{Theorem}
\newtheorem{proposition}[theorem]{Proposition}
\newtheorem{lemma}[theorem]{Lemma}
\newtheorem{corollary}[theorem]{Corollary}
\theoremstyle{definition}
\newtheorem{remark}[theorem]{Remark}
\theoremstyle{plain}

\newcommand{\val}[2]{[#1]_{#2}}
\newcommand{\NpD}{\mathcal{N}_{\mathrm{pD}}}
\newcommand{\DP}{\mathcal{D}_{\mathrm{P}}}
\newcommand{\Pset}{\mathcal{P}}
\newcommand{\seqnum}[1]{\href{https://oeis.org/#1}{\rm \underline{#1}}}

\title{Additive Representations by Primitive Dyck Numbers:\\ Exceptional Sets and Optimal Orders}
\author{Takayuki Kuriyama}
\date{}

\begin{document}

\maketitle

\begin{abstract}
Let $\NpD$ be the set consisting of $0$ and the integers represented in
binary by primitive Dyck words.  For a positive even integer $N$, let $r(N)$
be the least number of positive elements of $\NpD$ whose sum is $N$; the
associated sequence $a(n)=r(2n)$ is OEIS sequence \seqnum{A395858}.  We
determine the complete exceptional set for six-summand representability:
$46$ alone satisfies $r(46)=8$, the ten integers
$34,44,98,154,198,202,206,838,842,846$ satisfy $r(N)=7$, and every other
positive even integer satisfies $r(N)\leq6$.  Thus $848$ is the sharp even
threshold beyond which six summands always suffice.  We also determine the
exact relative asymptotic order.  For every $k\geq2$,
\[
  r(10\cdot4^{k+1}-6)=6,
\]
so six summands are necessary infinitely often and the asymptotic order is
exactly $6$.  The proof combines a regular sublanguage, finite interval
certificates, a self-similar five-summand propagation, and a base-$4$
characterization of the halved primitive family by two-coloured Motzkin words.
As a further consequence, that halved family is an asymptotic additive basis
of order exactly $5$.
\end{abstract}

\noindent\textbf{Keywords:} primitive Dyck number, additive
representation, additive basis, asymptotic order, exceptional set,
two-coloured Motzkin word, digital self-similarity, interval filling.

\noindent\textbf{2020 Mathematics Subject Classification:} Primary 11B13;
Secondary 05A15, 68Q45.

\section{Introduction}\label{sec:introduction}

Every string of balanced parentheses can be written over the alphabet
$\{1,0\}$, with $1$ for an opening bracket and $0$ for a closing one,
and so read as a binary integer.  Which integers arise this way, and how
additively rich is the set they form?  We take up a natural variant of
that question.

A \emph{Dyck word} is a word $w\in\{1,0\}^*$ with equally many $1$'s
and $0$'s in which every prefix has at least as many $1$'s as $0$'s;
thus $1$ opens a parenthesis and $0$ closes it.  A nonempty Dyck word is
\emph{primitive} if none of its proper nonempty prefixes is itself
balanced---equivalently, if the associated lattice path returns to
height zero for the first time only at its endpoint.  Such paths are
also commonly called \emph{irreducible} or \emph{indecomposable} Dyck
paths.  Let $\DP$ denote the language
of primitive Dyck words.  Ordered by length and then lexicographically,
its first eighteen elements are
\[
\begin{gathered}
10,\ 1100,\ 110100,\ 111000,\ 11010100,\ 11011000,\\
11100100,\ 11101000,\ 11110000,\ 1101010100,\ 1101011000,\\
1101100100,\ 1101101000,\ 1101110000,\ 1110010100,\\
1110011000,\ 1110100100,\ 1110101000.
\end{gathered}
\]

For a word $w$ over the base-$m$ digits, let $\val{w}{m}$ denote the
integer it represents in base $m$.  Adjoining $0$ for convenience, set
\[
 \NpD=\{0\}\cup\{\val{w}{2}:w\in\DP\}.
\]
The first eighteen positive elements are
\[
\begin{gathered}
2,12,52,56,212,216,228,232,240,\\
852,856,868,872,880,916,920,932,936;
\end{gathered}
\]
this is OEIS sequence \seqnum{A057547} \cite{OEIS-A057547}.  Every positive member of
$\NpD$ is even, since every nonempty Dyck word ends in $0$.  Thus only
even target integers can be represented by sums from $\NpD$.

The following sharper congruence observation will be used repeatedly.

\begin{lemma}\label{lem:mod4}
The number $2$ is the only positive primitive Dyck number congruent to
$2$ modulo $4$.  Every other positive primitive Dyck number is divisible
by $4$.
\end{lemma}

\begin{proof}
Every nonempty Dyck word ends in $0$.  A primitive Dyck word of length at
least four must in fact end in $00$: if it ended in $10$, the path would
already have returned to height zero just before its last two letters.
The word $10$ represents $2$, and every longer primitive word therefore
represents a multiple of $4$.
\end{proof}

Bell, Lidbetter, and Shallit studied the larger set obtained from
\emph{all} Dyck words, the \emph{totally balanced numbers} (OEIS sequence
\seqnum{A014486} \cite{OEIS-A014486}).  Using regular underapproximations and automata,
they proved that every even integer except
\[
        8,\ 18,\ 28,\ 38,\ 40,\ 82,\ 166
\]
is a sum of at most three totally balanced numbers, while two summands
do not suffice asymptotically \cite[Thm.\ 14]{BLS}.

The primitive restriction is not a cosmetic variant of that result.
Every nonempty Dyck word factors uniquely as a concatenation of
primitive Dyck words, whereas here each summand must consist of a single
factor.  Consequently, a representation by unrestricted Dyck numbers
does not in general yield a representation by the same number of
primitive Dyck numbers.

Digit-defined additive bases have also been studied by automata-theoretic
methods.  Bell, Hare, and Shallit characterized automatic sets that form
additive bases and gave an effective procedure for determining their
orders \cite{BHS}.  Related decision procedures have yielded sharp
additive results for binary palindromes \cite{RSS} and binary squares
\cite{MNRS}.  The language of primitive Dyck words is not regular, so
those general finite-automaton methods do not directly settle the
present problem.

Our proof differs at the structural point where regular approximation
alone ceases to suffice.  A regular sublanguage settles the congruence
class $0$ modulo $4$, while the class $2$ modulo $4$ is handled by a
self-similar closure property of the full primitive family after division
by $4$.  The same family admits a base-$4$ description by two-coloured
Motzkin words.  Its self-similarity has two complementary effects: it
propagates five-summand intervals, but its generation gaps also create an
infinite sequence of obstructions to four-summand representations.  These
facts yield the complete exceptional set, the optimal uniform bound, the
sharp eventual threshold, and the exact asymptotic order.

Within the Dyck words of length $2n$, the proportion that are primitive
tends to $1/4$ as $n\to\infty$.  Indeed, there are $C_n$ Dyck words of
length $2n$, and every primitive one is uniquely of the form $1u0$ with
$u$ a Dyck word of length $2n-2$; so there are $C_{n-1}$ primitive words
of length $2n$.  From the Catalan formula
$C_n=\frac1{n+1}\binom{2n}{n}$,
\[
 \frac{C_{n-1}}{C_n}=\frac{n+1}{2(2n-1)}\longrightarrow\frac14 .
\]
Its arithmetic effect, however, is striking.  Although all sufficiently
large even integers need only three unrestricted Dyck summands, the
primitive family has a small but nontrivial exceptional set.  We show
that precisely eleven positive even integers fail to be sums of at most
six primitive Dyck numbers.  Ten require seven summands, while $46$
alone requires eight.  All eleven exceptions lie below $848$.

The need for eight summands is entirely local: below $48$ the only
positive primitive Dyck numbers are $2$ and $12$.  The eventual
six-summand bound, on the other hand, comes from two self-similar
families.  A regular subfamily handles multiples of $4$, while the full
primitive family after division by $4$ handles the remaining congruence
class.

For a positive even integer $N$, let $r(N)$ be the least number of
positive primitive Dyck numbers whose sum is $N$, and define the central
integer sequence of this paper by
\[
        a(n)=r(2n),\qquad n\geq1.
\]
This sequence has been registered in the OEIS as
\seqnum{A395858} \cite{OEIS-A395858}.  Its first fifty terms are
\[
\begin{gathered}
1,2,3,4,5,1,2,3,4,5,6,2,3,4,5,6,7,3,4,5,6,7,8,4,5,\\
1,2,1,2,3,4,2,3,2,3,4,5,3,4,3,4,5,6,4,5,4,5,6,7,5.
\end{gathered}
\]
The following Python program reproduces the displayed terms and generates
50,000 terms of the sequence, making some recurring patterns visible:
\url{https://github.com/growupkuriyama-hub/lean_cfg_project/blob/master/LeanCfgProject/PrimDyck/A395858.py}.

\begin{theorem}\label{thm:classification}
The values of $r(N)$ exceeding six are exactly as follows:
\[
\begin{aligned}
 r(46)&=8,\\[4pt]
 r(N)=7\quad&\Longleftrightarrow\quad
 N\in\{34,44,98,154,198,202,206,838,842,846\}.
\end{aligned}
\]
Thus every other positive even integer satisfies $r(N)\leq6$.  In other
words, every even integer $N\geq848$ is a sum of at most six primitive
Dyck numbers.
\end{theorem}

Equivalently, for the OEIS sequence \seqnum{A395858}
\cite{OEIS-A395858}, defined by $a(n)=r(2n)$,
\[
        \max_{n\geq1}a(n)=8,
        \qquad
        a(n)=8\quad\Longleftrightarrow\quad n=23,
\]
and
\[
 a(n)=7
 \quad\Longleftrightarrow\quad
 n\in\{17,22,49,77,99,101,103,419,421,423\}.
\]
Thus the main theorem gives a complete determination of all terms of
$a(n)$ exceeding six, rather than only an eventual estimate.


The eventual bound six is also best possible.

\begin{theorem}\label{thm:asymptotic-order}
Let
\[
 h_{\mathrm{asym}}
 =\min\{h:\text{every sufficiently large even }N\text{ satisfies }r(N)\leq h\}.
\]
Then $h_{\mathrm{asym}}=6$.  More precisely, for every integer $k\geq2$,
\[
 N_k=10\cdot4^{k+1}-6
\]
satisfies $r(N_k)=6$.  Equivalently,
\[
 a(5\cdot4^{k+1}-3)=6\qquad(k\geq2).
\]
\end{theorem}

Thus the finite exceptional set above coexists with infinitely many large
integers for which the six-summand upper bound is attained.

The exceptional value $46$ already explains why the uniform bound can
never be smaller than eight.

\begin{proposition}\label{prop:46}
The integer $46$ requires exactly eight positive primitive Dyck numbers:
it is a sum of eight, but it is not a sum of at most seven.
\end{proposition}

\begin{proof}
Below $48$ the only positive primitive Dyck numbers are $2$ and $12$.
Indeed, the primitive words of length two and four are $10$ and $1100$,
whereas every primitive word of length at least six begins with $11$ and
therefore represents an integer at least $[110000]_2=48$.  Thus any
representation of $46$ has the form
\[
        46=12a+2b,\qquad a,b\geq0.
\]
Equivalently, $6a+b=23$.  Since $12\cdot4>46$, we have $a\leq3$, and
hence the number of summands satisfies
\[
        a+b=23-5a\geq8.
\]
Equality is attained by
\[
        46=12+12+12+2+2+2+2+2.
\]
\end{proof}

For example,
\[
 2026=2+240+852+932=2+216+872+936.
\]
The seven distinct integers appearing in these two representations all
have primitive Dyck binary expansions.

The sumset notation used below is standard in additive-basis theory
\cite[Chap.\ 1]{Nathanson}.  For a nonnegative integer $k$, let $kX$ denote the
set of sums of $k$ not necessarily distinct members of $X$, with
$0X=\{0\}$, and put
\[
 F_k(X)=\bigcup_{j=0}^{k}jX.
\]
Thus $F_k(X)$ is the set of sums of at most $k$ elements of $X$.  For a scalar
$c$, write $c\cdot X=\{cx:x\in X\}$ for the dilation of $X$.

The mod-$4$ property motivates the central halved family
\begin{equation}\label{eq:P-def}
 \Pset=\{v/4:v\in\NpD,\ v\geq12\},
\end{equation}
so that
\begin{equation}\label{eq:NpD-decomp}
 \NpD=\{0,2\}\cup4\Pset.
\end{equation}
The first elements of $\Pset$ are
\[
 3,13,14,53,54,57,58,60,213,214,217,\ldots.
\]


\section{The base-\texorpdfstring{$4$}{4} structure of
\texorpdfstring{$\Pset$}{P}}\label{sec:base4}

Assign weights to the base-$4$ digits by
\[
  \sigma(3)=1,\qquad \sigma(0)=-1,\qquad
  \sigma(1)=\sigma(2)=0.
\]
A word $w=c_1\cdots c_\ell\in\{0,1,2,3\}^{\ell}$ is called a
\emph{two-coloured Motzkin word} if every prefix has nonnegative total weight
and the total weight of $w$ is $0$.  The two zero-weight digits $1$ and $2$ are
regarded as distinct colours.  Write $[w]_4$ for the value of $w$ as a base-$4$
numeral, with $[\varepsilon]_4=0$.

\begin{lemma}[Base-$4$ characterization]\label{lem:base4}
A positive integer $p$ lies in $\Pset$ if and only if its base-$4$ expansion has
the form $3w$, where $w$ is a two-coloured Motzkin word.
\end{lemma}

\begin{proof}
Let $v=4p\in\NpD$ with $v\geq12$, and group the binary digits of $v$ into
consecutive pairs.  Reading the pairs as base-$4$ digits gives a word
$d_1\cdots d_L$.  At pair boundaries, the binary height is twice the cumulative
$\sigma$-weight, because
\[
  11\leftrightarrow3,\quad 10\leftrightarrow2,\quad
  01\leftrightarrow1,\quad 00\leftrightarrow0.
\]
Since $v$ is primitive and has length at least four, it begins with $11$ and
ends with $00$; hence $d_1=3$ and $d_L=0$.  At every proper pair boundary the
binary height is at least $2$, and at the final boundary it is $0$.

These boundary conditions also control heights inside each pair.  A pair of
type $01$ can lower the height by only one before restoring it; a nonfinal pair
of type $00$ ends at height at least $2$ and therefore starts at height at
least $4$; and the other two pair types begin with $1$.  Thus positivity before
the final bit is equivalent to positivity at all proper pair boundaries,
together with the final return to $0$.

Dividing $v$ by $4$ deletes the final base-$4$ digit $0$.  Hence the base-$4$
expansion of $p$ is $3w$, where every prefix of $3w$ has $\sigma$-weight at
least $1$ and the total weight of $3w$ is $1$.  Removing the initial digit $3$
shows that every prefix of $w$ has nonnegative weight and that $w$ has total
weight $0$.

Conversely, if $p$ has base-$4$ expansion $3w$ with $w$ a two-coloured Motzkin
word, append a final digit $0$ and expand each base-$4$ digit into its binary
pair.  The resulting binary word has equal numbers of $1$'s and $0$'s, remains
strictly above height $0$ before its final bit by the preceding within-pair
check, and returns to $0$ at the end.  It is therefore a primitive Dyck word,
so $4p\in\NpD$ and $p\in\Pset$.
\end{proof}

The characterization gives the self-similarity used later.

\begin{corollary}\label{cor:Pclosure}
If $p\in\Pset$, then $4p+1$ and $4p+2$ also belong to $\Pset$.
\end{corollary}

\begin{proof}
Write the base-$4$ expansion of $p$ as $3w$, with $w$ a two-coloured Motzkin
word.  Then $4p+1$ and $4p+2$ have expansions $3w1$ and $3w2$.  Appending either
zero-weight colour preserves the two-coloured Motzkin condition, so the claim
follows from Lemma~\ref{lem:base4}.
\end{proof}

We next determine the smallest and largest values in each generation.

\begin{lemma}[Extremal two-coloured Motzkin values]\label{lem:extremal}
For every $\ell\geq0$ and every two-coloured Motzkin word $w$ of length $\ell$,
\begin{equation}\label{eq:motzkin-extrema}
  \frac{4^{\ell}-1}{3}\leq[w]_4\leq4^{\ell}-2^{\ell}.
\end{equation}
Both bounds are attained.  The lower bound is attained by $1^{\ell}$, and the
upper bound is attained by
\[
  w_{\max,\ell}=
  \begin{cases}
    3^a0^a,&\ell=2a,\\
    3^a2\,0^a,&\ell=2a+1.
  \end{cases}
\]
\end{lemma}

\begin{proof}
We use induction on $\ell$.  The assertion is immediate for $\ell=0$.

\emph{Lower bound.}
The first digit cannot be $0$.  If $c_1\in\{1,2\}$, then the suffix
$w'=c_2\cdots c_\ell$ is again a two-coloured Motzkin word, and hence
\[
  [w]_4\geq4^{\ell-1}+[w']_4
  \geq4^{\ell-1}+\frac{4^{\ell-1}-1}{3}
  =\frac{4^{\ell}-1}{3}.
\]
If $c_1=3$, then
\[
  [w]_4\geq3\cdot4^{\ell-1}\geq\frac{4^{\ell}-1}{3}.
\]
Equality is attained by $w=1^{\ell}$.

\emph{Upper bound.}
If $c_1\in\{1,2\}$, then $w'$ is again a two-coloured Motzkin word, so
\[
  [w]_4
  \leq2\cdot4^{\ell-1}+\bigl(4^{\ell-1}-2^{\ell-1}\bigr)
  =3\cdot4^{\ell-1}-2^{\ell-1}
  \leq4^{\ell}-2^{\ell}.
\]

Now suppose $c_1=3$, and let $r$ be the least index at which the prefix weight
returns to $0$.  Then $2\leq r\leq\ell$, one has $c_r=0$, and
\[
  u=c_2\cdots c_{r-1},\qquad v=c_{r+1}\cdots c_\ell
\]
are two-coloured Motzkin words of lengths $r-2$ and $\ell-r$, respectively.
By place value and the induction hypothesis,
\begin{align*}
  [w]_4
  &=3\cdot4^{\ell-1}+[u]_4\,4^{\ell-r+1}+[v]_4\\
  &\leq3\cdot4^{\ell-1}
    +\bigl(4^{r-2}-2^{r-2}\bigr)4^{\ell-r+1}
    +4^{\ell-r}-2^{\ell-r}\\
  &=4^{\ell}-2^{2\ell-r}+2^{2\ell-2r}-2^{\ell-r}.
\end{align*}
It remains to verify
\[
  2^{2\ell-r}-2^{2\ell-2r}
  \geq2^{\ell}-2^{\ell-r}.
\]
After factoring, this is
\[
  2^{2\ell-2r}(2^r-1)\geq2^{\ell-r}(2^r-1),
\]
which follows from $r\leq\ell$.  This proves the upper bound.

Finally, the displayed words $w_{\max,\ell}$ are two-coloured Motzkin words.
A geometric-series calculation gives
\[
  [w_{\max,\ell}]_4=4^{\ell}-2^{\ell},
\]
so the upper bound is attained.
\end{proof}

An element of $\Pset$ belongs to \emph{generation $d$} if its base-$4$
expansion has exactly $d$ digits.  By Lemma~\ref{lem:base4}, a generation-$d$
element has value $3\cdot4^{d-1}+[w]_4$, where $w$ has length $d-1$.

\begin{corollary}[Generation bounds]\label{cor:gen}
For every $d\geq1$, every generation-$d$ element of $\Pset$ lies in
\begin{equation}\label{eq:generation-interval}
  [g_d,U_d]
  =\left[
  3\cdot4^{d-1}+\frac{4^{d-1}-1}{3},
  4^d-2^{d-1}
  \right],
\end{equation}
and both endpoints are attained.  In particular,
\begin{enumerate}[label=\textup{(\alph*)}]
\item there is no element of $\Pset$ in $(U_d,3\cdot4^d)$;
\item for every $k\geq0$,
\begin{equation}\label{eq:three-g}
  3g_{k+1}=10\cdot4^k-1.
\end{equation}
\end{enumerate}
\end{corollary}

\begin{proof}
The interval follows from Lemmas~\ref{lem:base4} and~\ref{lem:extremal}.  The
attainment statements follow from the extremal words in
Lemma~\ref{lem:extremal}.  Part~(a) follows because the next generation begins
at $g_{d+1}>3\cdot4^d$.  Finally,
\[
  3g_{k+1}
  =3\left(3\cdot4^k+\frac{4^k-1}{3}\right)
  =10\cdot4^k-1.
\]
\end{proof}

\section{A regular family and six-term interval filling}\label{sec:regular}

Inside the base-$4$ family described in Section~\ref{sec:base4}, a simple regular pattern runs through the primitive Dyck words.  For
example, the binary expansion of $868$ factors as
\[
 \operatorname{bin}(868)=1101100100=11\,01\,10\,01\,00
 \in 11(01\mid10)^*00 .
\]
Pairing the binary digits turns the initial block $11$ into the base-$4$
digit $3$, the final block $00$ into $0$, and each middle block $01$ or
$10$ into $1$ or $2$.  This suggests that sums of a fixed number of such
base-$4$ integers might fill long intervals.  The mixed sumset lemma
below makes the idea precise.

Let
\[
 \mathcal E=\{0\}\cup
 \bigl\{\val{3\varepsilon_1\cdots\varepsilon_k}{4}:
        k\geq0,\ \varepsilon_i\in\{1,2\}\bigr\}.
\]
When $k=0$ the string $\varepsilon_1\cdots\varepsilon_k$ is empty, so
the corresponding element is $\val{3}{4}=3$.  Multiplying by $4$ appends
a base-$4$ zero.  Replacing each base-$4$ digit by its two binary digits,
\[
\begin{aligned}
 3_{(4)}&\longleftrightarrow11_{(2)},&
 1_{(4)}&\longleftrightarrow01_{(2)},\\
 2_{(4)}&\longleftrightarrow10_{(2)},&
 0_{(4)}&\longleftrightarrow00_{(2)}.
\end{aligned}
\]
we find that every positive element of $4\cdot\mathcal E$ has binary
expansion in $11(01\mid10)^*00$.  After the initial $11$ the path stands
at height $2$; each middle block $01$ or $10$ returns it to height $2$;
and the final $00$ brings it to height zero for the first time.  Thus
every positive element of $4\cdot\mathcal E$ is a primitive Dyck number,
and, since $0\in\NpD$,
\begin{equation}\label{eq:4E}
        4\cdot\mathcal E\subseteq\NpD.
\end{equation}

The crux is that six elements of $\mathcal E$ already cover every integer
from $25$ on.

\begin{proposition}\label{prop:E}
Every integer $n\geq25$ belongs to $6\mathcal E$.
\end{proposition}

We prove this by a finite-digit interval argument.  The role of the
auxiliary sets below is worth explaining first.  When an $m$-digit element
of $\mathcal E_{m}$ is split into its leading digit and its lower digits, the
lower $m-1$ digits form a base-$4$ correction whose digits lie in
$\{0,1\}$.  The set $\Delta_m$ records exactly these corrections.  When
some summands in a six-term sum use a new leading digit, the remaining
lower-digit problem becomes a mixed sum of copies of $\Delta_{m-1}$ and
$\mathcal E_{m-1}$.  The mixed sumsets $S_{t,m}$ are designed to keep
track of all these possibilities simultaneously.

Set $u_0=0$ and, for $m\geq1$,
\begin{equation}\label{eq:umdef}
 u_m=[\underbrace{11\cdots1}_{m\text{ digits}}]_4
    =1+4+\cdots+4^{m-1}=\frac{4^m-1}{3},
\end{equation}
so that
\begin{equation}\label{eq:urecursion}
\begin{aligned}
 4^{m-1}&=3u_{m-1}+1 &&(m\geq1),\\
 u_{m-1}&=4u_{m-2}+1 &&(m\geq2).
\end{aligned}
\end{equation}
Define the correction set
\[
 \Delta_m=\Bigl\{[\delta_{m-1}\cdots\delta_0]_4:\delta_j\in\{0,1\}\Bigr\}
    =\Bigl\{\textstyle\sum_{j=0}^{m-1}\delta_j4^j:\delta_j\in\{0,1\}\Bigr\},
\]
so that $\max \Delta_m=u_m$.  Let $\mathcal E_{m}$ be the finite subset of
$\mathcal E$ consisting of $0$ together with the elements having at most
$m$ base-$4$ digits.  For example,
\[
\begin{aligned}
 \Delta_2&=\{[00]_4,[01]_4,[10]_4,[11]_4\}=\{0,1,4,5\},\\
 \mathcal E_2&=\{0,[3]_4,[31]_4,[32]_4\}=\{0,3,13,14\}.
\end{aligned}
\]
For $0\leq t\leq6$ set
\begin{equation}\label{eq:S-t-m}
 S_{t,m}=t\Delta_m+(6-t)\mathcal E_{m}.
\end{equation}
The largest element of $\mathcal E_{m}$ is
\[
 \max\mathcal E_{m}=[3\underbrace{22\cdots2}_{m-1\text{ digits}}]_4
                =3\cdot4^{m-1}+2u_{m-1},
\]
and for $m\geq2$ equation \eqref{eq:urecursion} gives
\begin{equation}\label{eq:maxErecursion}
 \max\mathcal E_{m-1}=11u_{m-2}+3=\frac{11u_{m-1}+1}{4}.
\end{equation}
Since $\max \Delta_m=u_m$,
\begin{equation}\label{eq:maxStm}
 \max S_{t,m}=tu_m+(6-t)\max\mathcal E_{m} .
\end{equation}
Moreover $\max\mathcal E_{m}>u_m$, so $\max S_{t,m}$ decreases with $t$.
Throughout, $[a,b]$ denotes the integer interval $\{a,a+1,\ldots,b\}$.

\begin{lemma}\label{lem:mixed}
For every $m\geq2$ and every $2\leq t\leq6$,
\begin{equation}\label{eq:mixedtge}
 S_{t,m}=[0,\max S_{t,m}],
\end{equation}
whereas
\begin{equation}\label{eq:mixedtone}
 S_{1,m}=[0,\max S_{1,m}]\setminus\{2\}.
\end{equation}
Moreover,
\begin{equation}\label{eq:mixedEminterval}
 [25,\ell_{m}]\subseteq6\mathcal E_{m},\qquad \ell_{m}=\frac{33}{8}\,4^m-4 .
\end{equation}
\end{lemma}

\begin{proof}
We first prove \eqref{eq:mixedtge} and \eqref{eq:mixedtone} in the base
case $m=2$.  Recall that
\[
 \Delta_2=\{0,1,4,5\},\qquad \mathcal E_2=\{0,3,13,14\}.
\]
Direct calculation gives $2\Delta_2=\{0,1,2,4,5,6,8,9,10\}$, and adding one
more copy of $\Delta_2$ fills every remaining gap, so $3\Delta_2=[0,15]$ and hence
\[
 t\Delta_2=[0,5t]\qquad(3\leq t\leq6).
\]
On the other hand,
\[
\begin{aligned}
2\mathcal E_2
={}&\{0,3,6,13,14,16,17,26,27,28\},\\[1ex]
3\mathcal E_2
={}&\{0,3,6,9\}\cup[13,14]\cup[16,17]\cup[19,20]\\
&\qquad\cup[26,31]\cup[39,42],\\[1ex]
4\mathcal E_2
={}&\{0,3,6,9\}\cup[12,14]\cup[16,17]\cup[19,20]\\
&\qquad\cup[22,23]\cup[26,34]\cup[39,45]\cup[52,56],\\[1ex]
5\mathcal E_2
={}&\{0,3,6,9\}\cup[12,17]\cup[19,20]\cup[22,23]\\
&\qquad\cup[25,37]\cup[39,48]\cup[52,59]\cup[65,70].
\end{aligned}
\]
Adding $\Delta_2$ to $5\mathcal E_2$, and $2\Delta_2$ to $4\mathcal E_2$, gives
\[
 S_{1,2}=\{0,1\}\cup[3,75],\qquad S_{2,2}=[0,66].
\]
Similarly,
\[
\begin{aligned}
S_{3,2}&=3\Delta_2+3\mathcal E_2=[0,15]+3\mathcal E_2=[0,57],\\[1ex]
S_{4,2}&=4\Delta_2+2\mathcal E_2=[0,20]+2\mathcal E_2=[0,48],\\[1ex]
S_{5,2}&=5\Delta_2+\mathcal E_2 =[0,25]+\mathcal E_2 =[0,39],\\[1ex]
S_{6,2}&=6\Delta_2=[0,30].
\end{aligned}
\]
Collecting these,
\[
\begin{array}{c|c}
t&S_{t,2}\\ \hline
1&[0,75]\setminus\{2\}\\
2&[0,66]\\
3&[0,57]\\
4&[0,48]\\
5&[0,39]\\
6&[0,30]
\end{array}
\]
which proves \eqref{eq:mixedtge} and \eqref{eq:mixedtone} for $m=2$.
For \eqref{eq:mixedEminterval} we must check
$[25,\ell_2]=[25,62]\subseteq6\mathcal E_2$.  The intervals $[26,28]$,
$[39,42]$, and $[52,56]$ are the sums of two, three, and four members of
$\{13,14\}$, respectively, so
\[
\begin{aligned}
25&=13+3+3+3+3+0,\\
[26,40]&=\{26,27,28\}+3\cdot\{0,1,2,3,4\},\\
[39,51]&=[39,42]+3\cdot\{0,1,2,3\},\\
[52,62]&=[52,56]+3\cdot\{0,1,2\}.
\end{aligned}
\]
This completes the base case.  The lower endpoint is sharp:
$24\notin6\mathcal E$, since the only positive elements of $\mathcal E$
below $25$ are $3,13,$ and $14$, and no sum of at most six of these is
$24$.

Now let $m\geq3$ and assume that
\eqref{eq:mixedtge}--\eqref{eq:mixedEminterval} hold at level $m-1$.
To lighten the notation during this step, set
\[
 u=u_{m-1},\qquad
 q=q_m=[3\underbrace{11\cdots1}_{m-1\text{ digits}}]_4
      =3\cdot4^{m-1}+u,\qquad
 M=\max\mathcal E_{m-1}.
\]
Then
\[
 4^{m-1}=3u+1,\qquad
 M=\frac{11u+1}{4}.
\]
Moreover,
\begin{equation}\label{eq:qmdef}
 q=10u+3.
\end{equation}
We shall repeatedly use
\begin{equation}\label{eq:qoverlap}
 2\cdot4^{m-1}+4u=q-1,
\end{equation}
whose two sides are both equal to $10u+2$.  Splitting according to the
leading base-$4$ digit gives
\[
 \Delta_m=\Delta_{m-1}\cup(4^{m-1}+\Delta_{m-1}),\qquad
 \mathcal E_{m}=\mathcal E_{m-1}\cup(q+\Delta_{m-1}).
\]
Indeed, every new $m$-digit element of $\Delta_m$ has leading digit $1$,
whereas a new $m$-digit element of $\mathcal E_{m}$ has leading digit $3$
followed by a lower $(m-1)$-digit correction from $\Delta_{m-1}$.  If $s$ of
the $t$ copies of $\Delta_m$ use their shifted part and $v$ of the $6-t$
copies of $\mathcal E_{m}$ use their new $m$-digit part, then the sumset in
\eqref{eq:S-t-m} decomposes as
\begin{equation}\label{eq:decomp}
 S_{t,m}=\bigcup_{v=0}^{6-t}\ \bigcup_{s=0}^{t}
 \bigl(s\cdot4^{m-1}+vq+S_{t+v,m-1}\bigr).
\end{equation}
Thus, for fixed $v$, the parameter $s$ produces $t+1$ translates by
multiples of $4^{m-1}$; increasing $v$ then moves to the next large block:
\[
 \underbrace{[vq,\ \cdots]}_{v\text{-block}}
 \qquad
 \underbrace{[(v+1)q,\ \cdots]}_{(v+1)\text{-block}}.
\]
The rest of the proof checks that neighboring translates and neighboring
$v$-blocks overlap.

We first treat $t\geq2$.  The induction hypothesis
\eqref{eq:mixedtge} gives $S_{t+v,m-1}=[0,\max S_{t+v,m-1}]$ for every
value of $v$ in \eqref{eq:decomp}.  For fixed $v$, consecutive values of
$s$ shift this interval by $4^{m-1}$.  Since $\max S_{j,m-1}$ decreases
with $j$,
\[
 \max S_{t+v,m-1}\geq\max S_{6,m-1}=6u
     \geq3u=4^{m-1}-1,
\]
so these $t+1$ translates overlap and form
\[
 [vq,\,vq+t\cdot4^{m-1}+\max S_{t+v,m-1}].
\]
The blocks for consecutive values of $v$ also overlap.  Whenever a next
block exists we have $t+v\leq5$, and since $t\geq2$,
\[
\begin{aligned}
 t\cdot4^{m-1}+\max S_{t+v,m-1}
 &\geq2\cdot4^{m-1}+\max S_{5,m-1}\\
 &=2\cdot4^{m-1}+5u+M\\
 &=(2\cdot4^{m-1}+4u)+(u+M)\\
 &\geq q-1,
\end{aligned}
\]
by \eqref{eq:qoverlap}.  Hence the union in \eqref{eq:decomp} is the
full integer interval from $0$ to its largest element $\max S_{t,m}$.
This proves \eqref{eq:mixedtge} at level $m$.

For $t=1$ the first block retains the single missing integer $2$.  For
each fixed $v\geq1$, the two $s$-translates involve $S_{1+v,m-1}$ with
$1+v\geq2$; they overlap because
\[
 \max S_{1+v,m-1}\geq\max S_{6,m-1}=6u\geq4^{m-1}-1.
\]
Every later $v$-block is therefore an interval.  Since $\max S_{j,m-1}$
decreases with $j$, it suffices to check the last pair of consecutive
$v$-blocks.  We have
\[
\begin{aligned}
 4^{m-1}+\max S_{5,m-1}-(q-1)
 &=M-2u-1\\
 &=\frac{11u+1}{4}-2u-1
  =\frac{3u-3}{4}\geq0,
\end{aligned}
\]
so consecutive $v$-blocks overlap.  Finally,
\[
\begin{aligned}
 \max S_{1,m-1}-(4^{m-1}+2)
 &=u+5M-4^{m-1}-2\\
 &=\frac{47u-7}{4}\geq0.
\end{aligned}
\]
Thus, within the first $v$-block, the unshifted long interval reaches
$4^{m-1}+2$ and joins the next $s$-translate without a gap.  Hence
\[
 S_{1,m}=[0,\max S_{1,m}]\setminus\{2\},
\]
which proves \eqref{eq:mixedtone} at level $m$.

It remains to propagate the interval inside $6\mathcal E_{m}$.  Taking
$t=0$ in \eqref{eq:decomp},
\[
 6\mathcal E_{m}=\bigcup_{v=0}^{6}\bigl(vq+S_{v,m-1}\bigr).
\]
The $v=0$ block is $S_{0,m-1}=6\mathcal E_{m-1}$, which contains
$[25,\ell_{m-1}]$ by the induction hypothesis.  Since $m\geq3$ we have
$u\geq5$, and
\[
 \ell_{m-1}-(q+2)=\frac{19u-39}{8}\geq0 .
\]
The $v=1$ block, namely $q+S_{1,m-1}$, contains
$\{q,q+1\}\cup[q+3,q+\max S_{1,m-1}]$ by the induction hypothesis.
Because the preceding interval reaches through $q+2$, the two join
without a gap.  To join the $v=j$ block to the $v=j+1$ block for
$j=1,2,3$, it is enough to have
\[
 \max S_{j,m-1}\geq q-1.
\]
Since $\max S_{j,m-1}$ decreases with $j$, the most restrictive case is
$j=3$; hence it suffices to check
\[
\begin{aligned}
 \max S_{3,m-1}-(q-1)
 &=3u+3M-(10u+2)\\
 &=\frac{5u-5}{4}\geq0 .
\end{aligned}
\]
Thus $\max S_{1,m-1},\ \max S_{2,m-1},\ \max S_{3,m-1}\geq q-1$, and
the blocks $v=1,2,3,4$ join in turn and cover through
\[
 4q+\max S_{4,m-1}=\frac{33}{8}\,4^m-4=\ell_{m} .
\]
The remaining blocks $v=5,6$ may begin after a gap, but they are not
needed for the interval propagated here.  This establishes
\eqref{eq:mixedtge}--\eqref{eq:mixedEminterval} at level $m$ and
completes the induction.
\end{proof}

Since $\ell_{m}\to\infty$ and $\mathcal E=\bigcup_{m\geq1}\mathcal E_{m}$,
Proposition \ref{prop:E} follows at once.

The preceding proposition handles the congruence class $0$ modulo
$4$: every multiple of $4$ at least $100$ is a sum of at most six
primitive Dyck numbers.

\section{A five-summand interval for the full halved family}
\label{sec:five-interval}

The regular family $\mathcal E$ is sufficient for multiples of $4$, but it
does not by itself provide the five-summand interval needed for the other
congruence class.  We now use the full family $\Pset$ from
\eqref{eq:P-def}, together with its self-similar closure in
Corollary~\ref{cor:Pclosure}.

We next establish a five-summand interval for $\Pset$.  The finite base
uses the two subsets
\begin{equation}\label{eq:L-def}
 L=\{3,13,14,53,54,57,58,60\}\subseteq\Pset.
\end{equation}
and
\begin{equation}\label{eq:H-def}
\begin{split}
 H=\{&213,214,217,218,220,229,230,233,234,236,\\
     &241,242,244,248\}\subseteq\Pset.
\end{split}
\end{equation}
The following tables display all the corresponding primitive Dyck
numbers.  In each entry the binary word $[4h]_2$ is primitive, and hence
$h\in\Pset$.
\[
\begin{array}{c|c@{\qquad}c|c}
\multicolumn{4}{c}{4L}\\
h&[4h]_2&h&[4h]_2\\ \hline
3&1100&53&11010100\\
13&110100&54&11011000\\
14&111000&57&11100100\\
58&11101000&60&11110000
\end{array}
\]
\[
\begin{array}{c|c@{\qquad}c|c}
\multicolumn{4}{c}{4H}\\
h&[4h]_2&h&[4h]_2\\ \hline
213&1101010100&233&1110100100\\
214&1101011000&234&1110101000\\
217&1101100100&236&1110110000\\
218&1101101000&241&1111000100\\
220&1101110000&242&1111001000\\
229&1110010100&244&1111010000\\
230&1110011000&248&1111100000
\end{array}
\]

The finite computations can be reproduced using the Python program
\path{verify_certificates_for_paper.py}, available as a
\href{https://github.com/growupkuriyama-hub/lean_cfg_project/blob/bc0211ccb478d79969f8ce6b5fe465e5f3fd9cf5/LeanCfgProject/PrimDyck/verify_certificates_for_paper.py}{fixed repository copy at commit \texttt{bc0211ccb478}}.
The same program is included as a supplementary file, together with a
README giving the execution command and the expected output.

The five statements in Lemma~\ref{lem:finitefive} are inclusion
certificates: the program checks that each displayed interval is
contained in the corresponding sumset, and does not claim that the
sumset equals that interval.  By contrast, for
Lemma~\ref{lem:finiteexception} the program computes the relevant sets
truncated to $[0,211]$ exactly and checks equality with the displayed
unions of intervals.

\begin{lemma}[Finite interval certificate]\label{lem:finitefive}
The following inclusions hold:
\begin{equation}\label{eq:finite-interval-certificate}
\begin{array}{c|c}
\text{sumset}&\text{integer interval contained}\\ \hline
5L &[215,254]\\
H+4L &[245,486]\\
2H+3L &[439,674]\\
3H+2L &[645,862]\\
4H+L &[855,1046].
\end{array}
\end{equation}
\end{lemma}

\begin{proof}
For the first inclusion in
\eqref{eq:finite-interval-certificate}, namely
$[215,254]\subseteq5L$, partition the set $L$ defined in
\eqref{eq:L-def} as
\[
 A=\{53,54,57,58,60\},\qquad B=\{3,13,14\},
\]
so that $L=A\sqcup B$.  Successive additions give
\[
 2A=[106,108]\cup[110,118]\cup\{120\},
\]
and
\[
 3A=[159,178]\cup\{180\},\qquad
 4A=[212,238]\cup\{240\}.
\]
Hence $4A+B=[215,254]$, proving the first inclusion in
\eqref{eq:finite-interval-certificate}.

For the second inclusion in
\eqref{eq:finite-interval-certificate}, namely
$[245,486]\subseteq H+4L$, repeated addition of $L$ gives
\begin{equation}\label{eq:4L-intervals}
\begin{split}
4L\supseteq{}&[32,34]\cup[42,45]\cup[52,56]\cup[82,102]\\
 &{}\cup[116,148]\cup[162,194]\cup[212,238].
\end{split}
\end{equation}
Among the seven intervals on the right-hand side of
\eqref{eq:4L-intervals}, adding the set $H$ defined in
\eqref{eq:H-def} to the first three short intervals gives
\[
\begin{aligned}
H+[32,34]&\supseteq[245,254]\cup[261,270]
 \cup[273,278]\cup[280,282],\\
H+[42,45]&\supseteq[255,265]\cup[271,281]\cup[283,293],\\
H+[52,56]&\supseteq[265,276]\cup[281,304].
\end{aligned}
\]
Their union contains $[245,304]$.  Since $\min H=213$,
$\max H=248$, and consecutive elements of $H$ differ by at most $9$,
if $I=[a,b]$ satisfies $b-a\geq8$, then
\[
 H+I=[213+a,248+b].
\]
Applying this to the final four intervals in
\eqref{eq:4L-intervals} gives, respectively,
\[
 [295,350],\quad[329,396],\quad[375,442],\quad[425,486].
\]
These intervals overlap, completing the proof that
$[245,486]\subseteq H+4L$.

For the third inclusion in
\eqref{eq:finite-interval-certificate}, namely
$[439,674]\subseteq2H+3L$, the following contained intervals suffice:
\[
 2H\supseteq[426,428]\cup[430,438]\cup[446,478]\cup\{496\},
\]
and
\[
\begin{split}
3L\supseteq{}&\{9\}\cup[19,20]\cup[39,42]\cup[73,77]
 \cup[109,111]\\
 &{}\cup[113,134]\cup[159,178]\cup\{180\}.
\end{split}
\]
The first two components of $2H$ give
\[
 [430,438]+9=[439,447],\quad
 [426,428]+[19,20]=[445,448],
\]
and
\[
 [430,438]+[19,20]=[449,458].
\]
Adding $[446,478]$ to, successively,
\[
\begin{gathered}
 \{9\},\ [39,42],\ [73,77],\ [109,111],\\
 [113,134],\ [159,178],\ \{180\}
\end{gathered}
\]
produces overlapping intervals covering $[455,658]$.  Finally,
\[
 496+[159,178]=[655,674].
\]
Together these intervals prove the third inclusion in
\eqref{eq:finite-interval-certificate}.

For the fourth inclusion in
\eqref{eq:finite-interval-certificate}, namely
$[645,862]\subseteq3H+2L$, use
\[
 3H\supseteq[639,734]\cup\{744\},\qquad
 2L\supseteq\{6\}\cup[70,74]\cup[110,118]\cup\{120\}.
\]
The four interval sums
\[
 [645,740],\quad[709,808],\quad[759,854],\quad[854,862]
\]
cover $[645,862]$.

For the fifth inclusion in
\eqref{eq:finite-interval-certificate}, namely
$[855,1046]\subseteq4H+L$, use
\[
 4H\supseteq[852,982]\cup[984,986],
 \qquad \{3,57,58,60\}\subseteq L.
\]
The intervals
\[
 [855,985],\quad[912,1042],\quad[1041,1044],\quad[1044,1046]
\]
then cover $[855,1046]$.
\end{proof}

Since the five intervals in Lemma~\ref{lem:finitefive} overlap, they give
\[
 [215,1046]\subseteq5\Pset.
\]

\begin{proposition}\label{prop:fiveP}
Every integer $n\geq215$ is a sum of exactly five elements of $\Pset$.
Consequently,
\[
 [212,\infty)\subseteq F_5(\Pset).
\]
The lower endpoint is sharp: none of the integers $209$, $210$, and
$211$ belongs to $F_5(\Pset)$.
\end{proposition}

\begin{proof}
Lemma~\ref{lem:finitefive} gives the stronger inclusion
$[215,1046]\subseteq5\Pset$; only the initial interval $[215,867]$ is
needed as the base case for strong induction on $n$.  Let $n\geq868$, and
choose the unique $\rho\in\{5,6,7,8\}$ such that
$\rho\equiv n\pmod4$.  Then
\[
 M=\frac{n-\rho}{4}\geq\frac{868-8}{4}=215.
\]
Since $M<n$, the induction hypothesis gives
\[
 M=q_1+\cdots+q_5,\qquad q_i\in\Pset.
\]
Because $5\leq\rho\leq8$, we may choose $\eta_i\in\{1,2\}$ such that
\[
 \eta_1+\cdots+\eta_5=\rho.
\]
Therefore
\[
 n=4M+\rho=\sum_{i=1}^{5}(4q_i+\eta_i).
\]
Corollary~\ref{cor:Pclosure} shows that every summand $4q_i+\eta_i$ belongs
to $\Pset$.  Thus $n\in5\Pset$, completing the induction.

Finally,
\[
\begin{aligned}
212&=53+53+53+53,\\
213&=53+53+53+54,\\
214&=53+53+54+54,
\end{aligned}
\]
so $[212,\infty)\subseteq F_5(\Pset)$.  We next verify that the only elements of $\Pset$ below $212$ are the
eight members of $L$.  Primitive Dyck words of lengths $4$, $6$, and $8$
give, after division by $4$, respectively,
\[
 \{3\},\qquad \{13,14\},\qquad \{53,54,57,58,60\}.
\]
Every primitive Dyck word of length $10$ has the form $1u0$, where $u$
is a Dyck word of length $8$.  The lexicographically least such word is
$1101010100$, whose value is $852$; after division by $4$ this gives
$213$.  Longer primitive words have still larger values.  Hence, with $L$ as
in \eqref{eq:L-def},
\[
 \Pset\cap[0,211]=L.
\]

Let
\[
 D=\{53,54,57,58,60\}\subseteq L.
\]
A sum of at most five members of $L$ using at most three terms from $D$
is at most
\[
 3\cdot60+2\cdot14=208,
\]
whereas a sum using at least four terms from $D$ is at least
\[
 4\cdot53=212.
\]
Hence $209,210,211\notin F_5(\Pset)$.
\end{proof}

\section{The complete exceptional set}\label{sec:exceptions}

We first translate six-summand representability in the congruence class
$2$ modulo $4$ into a finite-sum condition on $\Pset$.

\begin{lemma}\label{lem:sixcriterion}
Let $N=4m+2$.  Then $N$ is a sum of at most six positive primitive Dyck
numbers if and only if
\[
 m\in F_5(\Pset)
 \cup\bigl(1+F_3(\Pset)\bigr)
 \cup\bigl(2+F_1(\Pset)\bigr).
\]
\end{lemma}

\begin{proof}
By Lemma~\ref{lem:mod4}, every primitive Dyck summand other than $2$ is
a multiple of $4$.  Hence a representation of $N\equiv2\pmod 4$ contains
an odd number of copies of $2$.  If it has at most six summands, that
number is $1$, $3$, or $5$.

With one copy of $2$, division of the remaining sum by $4$ gives
$m\in F_5(\Pset)$.  With three copies it gives
$m-1\in F_3(\Pset)$, and with five copies it gives
$m-2\in F_1(\Pset)$.  This proves necessity.  Conversely, each of the
three memberships gives a representation of $N$ after multiplying the
corresponding $\Pset$-summands by $4$ and adjoining, respectively, one,
three, or five copies of $2$.  Thus the condition is also sufficient.
\end{proof}

\begin{lemma}[Small multiples of four]\label{lem:smallmultiples}
Among the positive multiples of $4$ below $100$, the only integer not
representable by at most six positive primitive Dyck numbers is $44$.
Moreover,
\[
 44=12+12+12+2+2+2+2.
\]
\end{lemma}

\begin{proof}
The positive primitive Dyck numbers below $100$ are precisely
\[
 2,\ 12,\ 52,\ 56.
\]
First suppose $N<52$ and write $N=4x$.  A sum of copies of $12$ and
pairs $2+2$ has the form
\[
 x=3a+c
\]
and uses $a+2c$ summands.  Under the condition $a+2c\leq6$, the possible
positive values of $x<13$ are
\[
 1,2,3,4,5,6,7,8,9,10,12;
\]
the sole missing value is $x=11$, corresponding to $N=44$.

For $52\leq N<100$, one first uses either $52=4\cdot13$ or
$56=4\cdot14$.  With at most five remaining summands, the residual
quotients obtainable from $12$ and pairs $2+2$ are
\[
 R=\{3a+c:a,c\geq0,\ a+2c\leq5\}
   =\{0,1,2,3,4,5,6,7,9,10,12,15\}.
\]
A direct inspection of these twelve integers gives
\[
 [13,24]\subseteq(13+R)\cup(14+R),
\]
so every multiple of $4$ from $52$ through $96$ has the required
representation.  The displayed seven-term representation of $44$
finishes the proof.
\end{proof}

\begin{lemma}[Finite exceptional-set certificate]\label{lem:finiteexception}
Let
\[
 L=\{3,13,14,53,54,57,58,60\}.
\]
Then
\[
\begin{aligned}
F_5(L)\cap[0,211]
={}&\{0,3,6,9\}\cup[12,17]\cup[19,20]\cup[22,23]\\
 &{}\cup[25,37]\cup[39,48]\cup[52,208].
\end{aligned}
\]
Moreover,
\[
\begin{aligned}
[0,211]\cap\Bigl(F_5(L)&\cup(1+F_3(L))\cup(2+F_1(L))\Bigr)\\
={}&[0,7]\cup[9,10]\cup[12,23]\cup[25,37]\\
 &{}\cup[39,48]\cup[52,208].
\end{aligned}
\]
Consequently, the complement in $[0,211]$ is
\[
 \{8,11,24,38,49,50,51,209,210,211\}.
\]
\end{lemma}

\begin{proof}
Starting with $0L=\{0\}$ and using the recursion
$jL=(j-1)L+L$, successive addition of the eight elements of $L$ gives,
up to $211$,
\[
\begin{aligned}
F_5(L)\cap[0,211]={}&\{0,3,6,9\}\cup[12,17]\cup[19,20]\cup[22,23]\\
 &{}\cup[25,37]\cup[39,48]\cup[52,208],
\end{aligned}
\]
while
\[
\begin{aligned}
(1+F_3(L))\cap[0,211]
={}&\{1,4,7,10\}\cup[14,15]\cup[17,18]\cup[20,21]\\
 &{}\cup[27,32]\cup[40,43]\cup[54,55]\cup[57,62]\\
 &{}\cup[64,65]\cup[67,78]\cup[80,89]\cup[107,135]\\
 &{}\cup[160,179]\cup\{181\},
\end{aligned}
\]
and
\[
 (2+F_1(L))\cap[0,211]
 =\{2,5\}\cup[15,16]\cup[55,56]\cup[59,60]\cup\{62\}.
\]
Taking the union of these displayed interval lists gives the second
formula, and its omitted integers are exactly those in the final set.
Each step involves only translating the preceding interval list by the
eight elements of $L$, so the calculation can be checked without an
exhaustive search.
\end{proof}

\begin{proof}[Proof of Theorem~\ref{thm:classification}]
Suppose first that $N\equiv0\pmod 4$.  If $N\geq100$, write $N=4n$.
Then $n\geq25$, so Proposition~\ref{prop:E} and \eqref{eq:4E} give a
representation of $N$ with at most six primitive Dyck summands.
Lemma~\ref{lem:smallmultiples} treats the remaining positive multiples
of $4$ and shows that $r(44)=7$.

Now let $N\equiv2\pmod 4$, say $N=4m+2$.  If $N\geq850$, then $m\geq212$,
so Proposition~\ref{prop:fiveP} gives $m\in F_5(\Pset)$.  By
Lemma~\ref{lem:sixcriterion}, $N$ is a sum of at most six primitive
Dyck numbers.  Together with the preceding paragraph, this already
shows that every even $N\geq848$ has such a representation.

It remains to classify the numbers $N\leq846$ with $N\equiv2\pmod 4$.
Here $0\leq m\leq211$, and Proposition~\ref{prop:fiveP} shows that the
only elements of $\Pset$ not exceeding $m$ lie in
\[
 L=\{3,13,14,53,54,57,58,60\}.
\]
Lemma~\ref{lem:finiteexception} shows that the values of $m$ not
covered by the criterion are
\[
 \{8,11,24,38,49,50,51,209,210,211\}.
\]
By Lemma~\ref{lem:sixcriterion}, the corresponding values $N=4m+2$
are precisely
\[
 34,46,98,154,198,202,206,838,842,846.
\]
Thus these ten numbers, together with $44$, are exactly the positive
even integers that are not sums of at most six primitive Dyck numbers.

All of them except $46$ have seven-summand representations:
\[
\begin{aligned}
34&=12+12+2+2+2+2+2,\\
44&=12+12+12+2+2+2+2,\\
98&=56+12+12+12+2+2+2,\\
154&=52+52+12+12+12+12+2,\\
198&=56+52+52+12+12+12+2,\\
202&=56+56+52+12+12+12+2,\\
206&=56+56+56+12+12+12+2,\\
838&=240+240+240+52+52+12+2,\\
842&=240+240+240+56+52+12+2,\\
846&=240+240+240+56+56+12+2.
\end{aligned}
\]
Proposition~\ref{prop:46} shows that $46$ requires exactly eight
summands.  This proves all assertions of the theorem.
\end{proof}

\begin{corollary}\label{cor:threshold}
The integer $848$ is the least even integer $T$ such that every even
integer $N\geq T$ is a sum of at most six primitive Dyck numbers.
\end{corollary}

\begin{proof}
Theorem~\ref{thm:classification} shows that six summands suffice for every even
integer at least $848$, whereas $846$ requires seven summands.
\end{proof}

\begin{corollary}\label{cor:eight}
Every nonnegative even integer is a sum of at most eight elements of
$\NpD$, and $46$ is the unique even integer that is not a sum of at most
seven elements of $\NpD$.
\end{corollary}

\begin{proof}
The assertion for positive integers follows immediately from
Theorem~\ref{thm:classification}; the integer $0$ is the empty sum.
\end{proof}


For $A\subseteq2\mathbb N_0$, we call $A$ a \emph{relative additive basis of
order at most $h$ for $2\mathbb N_0$} if every element of $2\mathbb N_0$ is a
sum of at most $h$ elements of $A$.  Its relative order is \emph{exactly $h$}
if $h$ is the least integer with this property.  We similarly call $A$ a
\emph{relative asymptotic additive basis of order at most $h$} if every
sufficiently large element of $2\mathbb N_0$ is a sum of at most $h$ elements
of $A$.

Theorem~\ref{thm:classification} gives the complete exceptional set for
six-summand representability.  Corollary~\ref{cor:eight} shows that $\NpD$ is
a relative additive basis of exact order eight for $2\mathbb N_0$, while
Corollary~\ref{cor:threshold} identifies the sharp six-summand threshold.
The next section proves that the relative asymptotic order is exactly six.


\section{Recurring gaps and asymptotic optimality}\label{sec:asymptotic}

Representations of an integer congruent to $2$ modulo $4$ split according to
the number of copies of the exceptional summand $2$.  For five summands the
criterion is as follows.

\begin{lemma}[Five-summand criterion]\label{lem:fivecrit}
Let $N=4m+2$ with $m\geq3$.  Then $r(N)\leq5$ if and only if
\begin{equation}\label{eq:fivecrit}
  m\in F_4(\Pset)\cup\bigl(1+F_2(\Pset)\bigr).
\end{equation}
\end{lemma}

\begin{proof}
By Lemma~\ref{lem:mod4}, a representation of $N$ uses some number $s\geq0$ of
copies of $2$ and otherwise uses summands from $4\Pset$.  Reduction modulo $4$
gives $2s\equiv2\pmod4$, so $s$ is odd.  If at most five summands are used,
then $s\in\{1,3,5\}$.  Writing the remaining summands as $4p_i$ and dividing
$N=2s+\sum_i4p_i$ by $4$ after subtracting $2$ gives
\[
  m=\frac{s-1}{2}+\sum_i p_i.
\]
For $s=1$ this gives $m\in F_4(\Pset)$, and for $s=3$ it gives
$m\in1+F_2(\Pset)$.  If $s=5$, then $m=2$, excluded by $m\geq3$.  This proves
necessity.  Conversely, membership in the first set of~\eqref{eq:fivecrit}
gives a representation using one copy of $2$, and membership in the second
gives one using three copies of $2$.
\end{proof}

The following lemma isolates the lower-bound mechanism.

\begin{lemma}[Recurring obstruction]\label{lem:obstruction}
For every integer $k\geq2$, put
\[
  x_k=10\cdot4^k-2.
\]
Then
\begin{equation}\label{eq:obstruction}
  x_k\notin F_4(\Pset)
  \qquad\text{and}\qquad
  x_k-1\notin F_2(\Pset).
\end{equation}
\end{lemma}

\begin{proof}
Fix $k\geq2$, and abbreviate $x=x_k$.  The least element of generation $k+2$
is
\[
  g_{k+2}>3\cdot4^{k+1}=12\cdot4^k>x.
\]
Consequently, every element of $\Pset$ that is at most $x$ has generation at
most $k+1$.  By Corollary~\ref{cor:gen}, every such element is either
\emph{small}, of generation at most $k$ and hence at most
\[
  U_k=4^k-2^{k-1},
\]
or \emph{large}, of generation $k+1$ and hence between $g_{k+1}$ and
$U_{k+1}=4^{k+1}-2^k$.

First consider $x-1=10\cdot4^k-3$.  It is larger than $U_{k+1}$ and smaller
than $g_{k+2}$, so it does not itself belong to $\Pset$.  Moreover, any sum of
two eligible elements is at most
\[
  2U_{k+1}=8\cdot4^k-2^{k+1}<10\cdot4^k-3=x-1.
\]
Thus $x-1\notin F_2(\Pset)$.

Now suppose, for contradiction, that
\[
  x=q_1+\cdots+q_t,\qquad t\leq4,\qquad q_i\in\Pset.
\]
Let $\lambda$ be the number of large summands.  If $\lambda\leq2$, then filling
the remaining positions with the largest possible small summands gives
\begin{align*}
  x
  &\leq2U_{k+1}+2U_k\\
  &=2(4^{k+1}-2^k)+2(4^k-2^{k-1})\\
  &=10\cdot4^k-3\cdot2^k
  <10\cdot4^k-2=x,
\end{align*}
a contradiction.  Hence $\lambda\geq3$.  But then the three large summands
alone have sum at least
\[
  3g_{k+1}=10\cdot4^k-1=x+1
\]
by~\eqref{eq:three-g}, again a contradiction.  Therefore
$x\notin F_4(\Pset)$.
\end{proof}

\begin{proof}[Proof of Theorem~\ref{thm:asymptotic-order}]
Fix $k\geq2$ and put $m=x_k=10\cdot4^k-2$.  Then
\[
  N_k=4m+2=10\cdot4^{k+1}-6.
\]
By Lemma~\ref{lem:obstruction},
\[
  m\notin F_4(\Pset),\qquad m-1\notin F_2(\Pset).
\]
The five-summand criterion, Lemma~\ref{lem:fivecrit}, therefore gives
$r(N_k)\geq6$.  The integer $N_k$ is not one of the eleven exceptions in
Theorem~\ref{thm:classification}, so $r(N_k)\leq6$.  Hence $r(N_k)=6$.
Since $N_k\to\infty$, no $h\leq5$ bounds $r(N)$ for all sufficiently large even
$N$, while Theorem~\ref{thm:classification} gives the upper bound $h=6$.
Therefore $h_{\mathrm{asym}}=6$.
\end{proof}

\begin{corollary}[Order of $\Pset$]\label{cor:orderfive}
The set $\Pset$ is an asymptotic additive basis of order exactly $5$ for the
nonnegative integers.  More precisely, every integer $n\geq212$ lies in
$F_5(\Pset)$, whereas $F_4(\Pset)$ and $F_3(\Pset)$ each omit infinitely many
integers.
\end{corollary}

\begin{proof}
Proposition~\ref{prop:fiveP} gives $[212,\infty)\subseteq F_5(\Pset)$.  For
every $k\geq2$, Lemma~\ref{lem:obstruction} gives
\[
  x_k=10\cdot4^k-2\notin F_4(\Pset).
\]
Since $x_k\to\infty$, four summands do not suffice asymptotically.  The same is
true for three summands because $F_3(\Pset)\subseteq F_4(\Pset)$.
\end{proof}

\begin{corollary}[Multiples of $4$ requiring five summands]
\label{cor:five-multiples}
For every integer $k\geq3$,
\[
  M_k=10\cdot4^{k+1}-8
\]
satisfies $r(M_k)=5$.  Hence infinitely many multiples of $4$ require exactly
five positive primitive Dyck summands.
\end{corollary}

\begin{proof}
Write $M_k=4x_k$.  Since $k\geq3$, one has $x_k\geq212$, and
Proposition~\ref{prop:fiveP} gives $x_k\in F_5(\Pset)$.  Multiplying such a
representation by $4$ shows that $r(M_k)\leq5$.

Suppose that $M_k$ had a representation using at most four primitive Dyck
summands, and let $s$ be the number of copies of $2$.  Since
$M_k\equiv0\pmod4$, $s$ is even, so $s\in\{0,2,4\}$.  If $s=0$, division by
$4$ gives $x_k\in F_4(\Pset)$, contradicting Lemma~\ref{lem:obstruction}.  If
$s=2$, division after removing the two copies of $2$ gives
$x_k-1\in F_2(\Pset)$, again a contradiction.  If $s=4$, no other summand is
available and division by $4$ would give $x_k=2$, which is impossible.  Thus
$r(M_k)\geq5$, and equality follows.
\end{proof}

\begin{remark}[The two residue classes]
If $N=4n$, then $n\in F_5(\Pset)$ for all sufficiently large $n$, so
$r(N)\leq5$ eventually; Corollary~\ref{cor:five-multiples} shows that equality
occurs infinitely often.  If $N=4m+2$, the number of copies of $2$ must be odd.
For $m=x_k$, the one-copy case would require $m\in F_4(\Pset)$ and the
three-copy case would require $m-1\in F_2(\Pset)$; Lemma~\ref{lem:obstruction}
excludes both.  This residue class raises the relative asymptotic order from
$5$ to $6$.
\end{remark}

\section{OEIS sequences and computational verification}
\label{sec:oeis}

Three OEIS sequences occur naturally.  The totally balanced numbers form
\seqnum{A014486}~\cite{OEIS-A014486}; the positive primitive Dyck numbers form
\seqnum{A057547}~\cite{OEIS-A057547}; and the central sequence of this paper is
\[
  a(n)=r(2n),
\]
namely \seqnum{A395858}~\cite{OEIS-A395858}.  Theorem~\ref{thm:classification}
determines all terms exceeding six:
\[
  a(n)=8\iff n=23,
\]
and
\[
  a(n)=7\iff
  n\in\{17,22,49,77,99,101,103,419,421,423\}.
\]
Theorem~\ref{thm:asymptotic-order} supplies an explicit infinite subsequence at
level six:
\begin{equation}\label{eq:A395858-six-subsequence}
  a(5\cdot4^{k+1}-3)=6\qquad(k\geq2).
\end{equation}

The program 
\url{https://github.com/growupkuriyama-hub/lean_cfg_project/blob/master/LeanCfgProject/PrimDyck/A395858.py}
reproduces the displayed initial terms and generates the first $50{,}000$ terms
of A395858.  The supplementary program
\path{verify_certificates_for_paper.py} checks the five finite interval
inclusions, the sharp lower endpoint for Proposition~\ref{prop:fiveP}, the
finite exceptional-set certificate, and an independent finite-range dynamic
programming computation of $r(N)$.  These computations provide reproducible
cross-checks; the infinite assertions are proved theoretically above.

\section{Concluding remarks}

The results determine three exact orders associated with primitive Dyck
numbers.  The relative additive order on all nonnegative even integers is
$8$, the relative asymptotic order is $6$, and the asymptotic additive order of
$\Pset$ is $5$.  The same base-$4$ self-similarity that propagates long
five-summand intervals also produces the recurring obstruction
$x_k=10\cdot4^k-2$.

It remains to characterize the complete set of even integers $N$ for which
$r(N)=6$.  Computations suggest additional block structure near the points
$10\cdot4^k$, but the exact widths and endpoints of those blocks are not yet
proved.  A broader question is whether analogous regular-subfamily upper
bounds and generation-gap lower bounds occur for other indecomposable digital
languages.

\begin{thebibliography}{9}

\bibitem{BHS}
J.~P. Bell, K.~G. Hare, and J.~Shallit,
When is an automatic set an additive basis?,
\emph{Proc. Amer. Math. Soc. Ser. B} \textbf{5} (2018), 50--63.

\bibitem{BLS}
J.~P. Bell, T.~F. Lidbetter, and J.~Shallit,
Additive number theory via approximation by regular languages,
\emph{Internat. J. Found. Comput. Sci.} \textbf{31} (2020), 667--687.

\bibitem{MNRS}
P.~Madhusudan, D.~Nowotka, A.~Rajasekaran, and J.~Shallit,
Lagrange's theorem for binary squares,
in \emph{43rd International Symposium on Mathematical Foundations of
Computer Science}, Leibniz Int. Proc. Inform., Vol.~117, Schloss
Dagstuhl--Leibniz-Zentrum f{\"u}r Informatik, 2018, pp.~18:1--18:14.

\bibitem{Nathanson}
M.~B. Nathanson,
\emph{Additive Number Theory: The Classical Bases},
Grad. Texts in Math., Vol.~164, Springer, 1996.

\bibitem{OEIS-A014486}
OEIS Foundation Inc.,
\emph{The On-Line Encyclopedia of Integer Sequences, A014486},
\url{https://oeis.org/A014486}.

\bibitem{OEIS-A057547}
OEIS Foundation Inc.,
\emph{The On-Line Encyclopedia of Integer Sequences, A057547},
\url{https://oeis.org/A057547}.

\bibitem{OEIS-A395858}
OEIS Foundation Inc.,
\emph{The On-Line Encyclopedia of Integer Sequences, A395858},
\url{https://oeis.org/A395858}.

\bibitem{RSS}
A.~Rajasekaran, J.~Shallit, and T.~Smith,
Sums of palindromes: an approach via automata,
in \emph{35th Symposium on Theoretical Aspects of Computer Science},
Leibniz Int. Proc. Inform., Vol.~96, Schloss Dagstuhl--Leibniz-Zentrum
f{\"u}r Informatik, 2018, pp.~54:1--54:12.

\end{thebibliography}


\begin{flushleft}
Takayuki Kuriyama\\
Growup Learning Academy\\
4-8-3 Shakujii-machi\\
Nerima-ku, Tokyo 177-0041\\
Japan\\
\textit{Email address}: \href{mailto:growup.kuriyama@gmail.com}
{\texttt{growup.kuriyama@gmail.com}}
\end{flushleft}

\end{document}
